\subsection*{1.2. 集合层}

\subsubsection*{1.2.1.}
设 $D$ 是圆盘（1.1.6），$0\in D$，并令 $D^*=D-\{0\}$。用 $\mathcal D^*$ 表示 $D^*$ 上局部常值集合层的拓扑斯。由 1.1.8，若将 $D$ 换为包含 $0$ 的较小圆盘 $D'$，此拓扑斯不会“改变”；当 $D$ 带有可微流形结构时，将 $D$ 换为其在 $0$ 处的切空间也同样如此。

对 $t\in D^*$，函子 $F\longmapsto F_t$ 将此拓扑斯等同于带有 $\pi_1(D^*)$ 作用的集合的拓扑斯（1.1.2）。由 1.1.4，函子
\[
F\longmapsto
\bigl(\widetilde D^*\longmapsto
H^0(\widetilde D^*,p^*F)\bigr)
\]
将其等同于下述函子的拓扑斯
\[
(\text{万有覆盖，底空间为 }D^*)
\longrightarrow(\mathrm{Sets}),
\]
这些等同由 (1.1.4.1) 联系起来。

\subsubsection*{1.2.2.}
设 $F$ 是 $D$ 上的层，其在 $D^*$ 上的限制局部常值。存在唯一的层态射
\[
\underline{F_0}\xrightarrow{\ \alpha\ }F
\]
使其在 $0$ 处的茎上的映射为恒等映射。如 XIII, 0.1 所述（取 $X_0=(P^t)$），函子
\[
F\longmapsto(F_0,F|_{D^*},\alpha)
\]
将这些层等同于由一个集合 $F_0$、$D^*$ 上的一个局部常值层 $F_\eta$ 及一个态射
\[
\alpha:
(\text{取值为 }F_0\text{ 的常值层，底空间为 }D^*)
\longrightarrow F_\eta.
\]
构成的三元组。

\subsubsection*{1.2.3.}
用 $\mathcal D$ 表示 $D$ 上限制到 $D^*$ 后局部常值的集合层所构成的拓扑斯。$\mathcal D$ 到 $D$ 上层的拓扑斯的包含，是一个拓扑斯态射 $D\longrightarrow\mathcal D$ 的逆像函子，此态射嵌入交换图
\[
\begin{tikzcd}
\{0\} \arrow[r, hook] \arrow[d, equal]
  & D \arrow[d, "\beta"]
  & D^* \arrow[l, hook'] \arrow[d, "\beta"] & {} \\
\{0\} \arrow[r, hook]
  & \mathcal D
  & \mathcal D^* \arrow[l, hook'] & {}.
\end{tikzcd}
\]
如 1.1.1 中一样，可验证：将 $D$ 换为较小圆盘或 $D$ 在 $0$ 处的切空间时，此拓扑斯不变。

\subsubsection*{1.2.4.}
设 $Y$ 是拓扑空间。积拓扑斯 $Y\times\mathcal D$ 有以下三种描述（其等价性见 XIII, 0.1）。

\medskip
\noindent\textit{(a)} $Y\times\mathcal D$ 是 $Y\times D$ 上满足下列条件的层 $F$ 所构成的拓扑斯：

\smallskip
\noindent $(*)$ $D^*$ 可由开子集 $U$ 覆盖，使 $F|_{Y\times U}$ 是 $Y$ 上某个层的逆像。

\medskip
\noindent\textit{(b)} 设 $t\in D^*$。拓扑斯 $Y\times\mathcal D$ 是三元组 $(F_0,F_t,\alpha)$ 的拓扑斯，其中

\smallskip
\noindent $\alpha)$ $F_0$ 是 $Y$ 上的层；

\noindent $\beta)$ $F_t$ 是 $Y$ 上带有 $\pi_1(D^*)$ 作用的层；

\noindent $\gamma)$ $\alpha:F_0\longrightarrow F_t$ 是像包含于不变部分的态射。

\medskip
\noindent\textit{(c)} 拓扑斯 $Y\times\mathcal D$ 是三元组 $(F_0,F_\eta,\alpha)$ 的拓扑斯，其中

\smallskip
\noindent $\alpha')$ $F_0$ 是 $Y$ 上的层；

\noindent $\beta')$ $F_\eta$ 是下述函子，其值记为 $F_{\bar\eta}$：
\[
(\text{万有覆盖，底空间为 }D^*)
\longrightarrow(\text{层，底空间为 }Y);
\]

\noindent $\gamma')$ $\alpha$ 是从 $F_0$（视为从 $D^*$ 的万有覆盖到 $Y$ 上的层、取常值 $F_0$ 的函子）到 $F_\eta$ 的态射。

此处并不在意这个拓扑斯是否确实是 $Y$ 与 $\mathcal D$ 的积拓扑斯。同样，可以任选以下一种方式定义 $Y\times\mathcal D^*$：

\medskip
\noindent\textit{(a$^*$)} $Y\times D^*$ 上满足 $(*)$ 的层所构成的拓扑斯；

\noindent\textit{(b$^*$)} 对 $t\in D^*$，是 $Y$ 上带有 $\pi_1(D^*)$ 作用的层 $F_t$ 所构成的拓扑斯；

\noindent\textit{(c$^*$)} 下述函子的拓扑斯：
\[
(\text{万有覆盖，底空间为 }D^*)
\longrightarrow(\text{层，底空间为 }Y).
\]

